\chapter{Implicaciones de política pública}

\section{La consolidación es real y recae entera sobre el gasto}

El Paquete Económico 2025 reduce los requerimientos financieros del sector público
de 5.9 a 3.9\% del PIB y convierte un déficit primario de 1.4 puntos en un
superávit de 0.6. \textbf{Dos puntos del PIB de ajuste en un ejercicio, sin una
sola medida de ingreso.} No hay miscelánea fiscal: la palabra no aparece en la
iniciativa de Ley de Ingresos.

El ajuste recae sobre el gasto programable, que cae de 19.7 a 17.8\% del PIB. Y
dentro del gasto programable recae sobre la obra: la inversión física baja de 3.0
a 2.3\% del PIB y el capítulo de obra pública directa cae 37.7\% real, mientras
las inversiones financieras se mantienen.

\section{El acervo no cede, y una parte de por qué no cede es un supuesto}

El saldo histórico de los requerimientos financieros se queda en 51.4\% del PIB, y
el escenario oficial lo mantiene ahí hasta 2030. La descomposición del capítulo 12
muestra que los 3.9 puntos de requerimientos se compensan con 3.18 de crecimiento
del PIB nominal y 0.79 de apreciación del peso.

\textbf{Esas ocho décimas dependen de que el tipo de cambio cierre en 18.5 pesos
por dólar.} Si cerrara donde cerró 2024, la razón subiría a 52.2 sin que nada de
la política fiscal cambiara. Es la sensibilidad más grande del paquete y no está
declarada como tal.

Y la inflación no ayuda: descompuesta, resta 2.12 puntos por el denominador y el
sector público paga 2.07 de compensación sobre la deuda indexada. \textbf{El
efecto neto es de cinco centésimas.} La idea de que la inflación licúa la deuda no
se sostiene cuando la tasa nominal la sigue.

\section{Lo comprometido crece más rápido que lo que lo paga}

Los gastos obligatorios con pensiones equivalen al 82.6\% del gasto neto total y
al 145\% de la recaudación tributaria. Las pensiones, contributivas y no
contributivas, suman 5.99\% del PIB. El costo financiero sube 5.4\% real y es el
ramo que más crece del presupuesto.

\textbf{Las dos razones estructurales de deuda empeoran pese a la
consolidación}: el costo financiero pasa de 0.256 a 0.262 de los impuestos, y el
acervo de 3.40 a 3.51 años de recaudación. La consolidación es del flujo; el
acervo todavía no la registra.

\section{Dentro del gasto social, la reasignación es nítida}

Educación se mantiene plana y reasigna: la beca universal de educación básica sube
26,826.5 millones de pesos reales, mientras los servicios de media superior, el
subsidio a universidades estatales y La Escuela es Nuestra pierden entre 4,500 y
5,400 cada uno, y el programa de infraestructura educativa desaparece.
\textbf{Transferencia monetaria a las familias contra servicios e infraestructura.}

Salud cae, y no toda la caída es lo que parece: descontando un adeudo liquidado
que no se repite, el retroceso real pasa de 12.2 a 7.2\%, y su origen está
identificado en el fondo de aportaciones para los servicios de salud, que cae
42.6\% real. \textbf{Los institutos que prestan el servicio crecen; quien pierde
es la población sin seguridad social atendida por los estados.}

Seguridad cae 17.6\% real, concentrado en la Guardia Nacional, que pasa a menos de
la mitad sin cambiar de ramo ni de nombre. \textbf{El componente militar no cae};
lo que baja en Defensa es obra ferroviaria que se traslada a la administración
civil.

\section{Lo que el paquete no permite saber}

Cuatro cosas, y las cuatro se declaran en sus capítulos y no solo aquí.

\begin{enumerate}
\item \textbf{El gasto en salud por afiliado}, porque no hay población por
afiliación.
\item \textbf{El gasto federalizado por habitante}, porque no hay población por
entidad.
\item \textbf{El gasto educativo por alumno}, porque el paquete no declara
matrícula, aunque obliga a reportarla dos veces en el decreto.
\item \textbf{Qué parte del desplazamiento de la renta petrolera hacia el ingreso
propio de las empresas es política}, porque falta la trayectoria de la tasa del
derecho por la utilidad compartida.
\end{enumerate}

\section{La recomendación de esta casa}

Una sola, y es de transparencia antes que de política.

\textbf{El Paquete Económico 2025 se publicó sin capa de texto.} Noventa de las
noventa y una páginas de los Criterios Generales son imagen. Un documento que
fundamenta nueve billones de pesos de gasto y que no se puede buscar, citar ni
verificar por medios ordinarios impone un costo a cualquiera que quiera
discutirlo, y ese costo no es neutral: recae sobre quien tiene menos recursos para
absorberlo.

Este documento pudo trabajar con esas páginas porque las leyó como imágenes y
validó cada cifra con cincuenta pruebas de identidad contable. \textbf{No debería
haber hecho falta.} Publicar el mismo archivo con su texto seleccionable no cuesta
nada y devuelve la discusión presupuestaria a quien quiera participar en ella.

\clearpage
